\section{Conclusão}
\label{cap:conclusao}

A execução deste trabalho permitiu comprovar de forma objetiva e prática a profunda interdependência teórica existente entre os conceitos matemáticos abordados\textsuperscript{}. Através das simulações computacionais desenvolvidas e analisadas, evidenciou-se que essas ferramentas formam um ciclo contínuo: a Teoria da Probabilidade baseia o cenário inicial mapeando as chances, o Teorema de Bayes atualiza a nossa visão assim que novas evidências surgem, o Valor Esperado direciona a tomada de decisão ao calcular tendências a longo prazo, e a Variância alerta medindo a instabilidade e o risco real da operação\textsuperscript{}. Fica claro, portanto, que utilizar essas métricas de maneira isolada compromete a fidelidade da análise estatística\textsuperscript{}.

Diante dos resultados obtidos, conclui-se sobre a indiscutível importância prática de se dominar esses quatro fundamentos estatísticos em conjunto\textsuperscript{}. A compreensão integrada dessas métricas é absolutamente essencial para a tomada de decisões reais e seguras em qualquer projeto computacional ou empresarial que envolva cenários de incerteza\textsuperscript{}. Ao abandonar o empirismo e aplicar essa engrenagem metodológica, desenvolvedores e engenheiros tornam-se capazes de mensurar riscos de forma assertiva, evitando armadilhas matemáticas e garantindo a construção de sistemas lógicos, eficientes e robustos diante do imprevisível\textsuperscript{}.
